% Copyright (C) 2020 Saurabh Joshi
%% 
\let\negmedspace\undefined
\let\negthickspace\undefined

\documentclass[journal,12pt,twocolumn]{IEEEtran}
%\documentclass[journal,12pt,twocolumn]{IEEEtran}
%
\usepackage{setspace}
\usepackage{gensymb}
%\doublespacing
\singlespacing

%\usepackage{graphicx}
%\usepackage{amssymb}
%\usepackage{relsize}
\usepackage[cmex10]{amsmath}
%\usepackage{amsthm}
%\interdisplaylinepenalty=2500
%\savesymbol{iint}
%\usepackage{txfonts}
%\restoresymbol{TXF}{iint}
%\usepackage{wasysym}
\usepackage{amsthm}
\usepackage{mathrsfs}
\usepackage{txfonts}
\usepackage{stfloats}
\usepackage{cite}
\usepackage{cases}
\usepackage{subfig}
%\usepackage{xtab}
\usepackage{longtable}
\usepackage{multirow}
%\usepackage{algorithm}
%\usepackage{algpseudocode}
\usepackage{enumitem}
\usepackage{mathtools}
\usepackage{tikz}
\usepackage{circuitikz}
\usepackage{verbatim}
\usepackage{hyperref}
%\usepackage{stmaryrd}
\usepackage{tkz-euclide} % loads  TikZ and tkz-base
%\usetkzobj{all}
\usepackage{listings}
    \usepackage{color}                                            %%
    \usepackage{array}                                            %%
    \usepackage{longtable}                                        %%
    \usepackage{calc}                                             %%
    \usepackage{multirow}                                         %%
    \usepackage{hhline}                                           %%
    \usepackage{ifthen}                                           %%
  %optionally (for landscape tables embedded in another document): %%
    \usepackage{lscape}     
\usepackage{multicol}
\usepackage{chngcntr}
\usepackage{iftex}
%\usepackage[latin9]{inputenc}
\usepackage{geometry}
\usepackage{bm}
%\geometry{verbose,tmargin=2cm,bmargin=3cm,lmargin=1.8cm,rmargin=1.5cm,headheight=2cm,headsep=2cm,footskip=3cm}
\usepackage{array}
\newcolumntype{L}[1]{>{\raggedright\let\newline\\\arraybackslash\hspace{0pt}}m{#1}}
\newcolumntype{C}[1]{>{\centering\let\newline\\\arraybackslash\hspace{0pt}}m{#1}}
\newcolumntype{R}[1]{>{\raggedleft\let\newline\\\arraybackslash\hspace{0pt}}m{#1}}

%\usepackage{graphicx}
%\usepackage{setspace}
%\usepackage{parskip}

\def \hsp {\hspace{3mm}}

\makeatletter

\providecommand{\tabularnewline}{\\}



\makeatother
\ifxetex
\usepackage[T1]{fontenc}
\usepackage{fontspec}
%\setmainfont[ Path = fonts/]{Sanskrit_2003.ttf}
\newfontfamily\nakulafont[Script=Devanagari,AutoFakeBold=2,Path = fonts/]{Nakula}
%\newfontfamily\liberationfont{Liberation Sans Narrow}
%\newfontfamily\liberationsansfont{Liberation Sans}
\fi
\usepackage{tikz}
\usepackage{xcolor}
%\usepackage{enumerate}

%\usepackage{wasysym}
%\newcounter{MYtempeqncnt}
\DeclareMathOperator*{\Res}{Res}
%\renewcommand{\baselinestretch}{2}
\renewcommand\thesection{\arabic{section}}
\renewcommand\thesubsection{\thesection.\arabic{subsection}}
\renewcommand\thesubsubsection{\thesubsection.\arabic{subsubsection}}

\renewcommand\thesectiondis{\arabic{section}}
\renewcommand\thesubsectiondis{\thesectiondis.\arabic{subsection}}
\renewcommand\thesubsubsectiondis{\thesubsectiondis.\arabic{subsubsection}}
\makeatletter 
\renewcommand{\thefigure}{S\@arabic\c@figure}
\makeatother

% correct bad hyphenation here
\hyphenation{op-tical net-works semi-conduc-tor}
\def\inputGnumericTable{}                                 %%

\lstset{
language=tex,
frame=single, 
breaklines=true
}

%\begin{document}
%


\newtheorem{theorem}{Theorem}[section]
\newtheorem{problem}{Problem}
\newtheorem{proposition}{Proposition}[section]
\newtheorem{lemma}{Lemma}[section]
\newtheorem{corollary}[theorem]{Corollary}
\newtheorem{example}{Example}[section]
\newtheorem{definition}[problem]{Definition}
%\newtheorem{thm}{Theorem}[section] 
%\newtheorem{defn}[thm]{Definition}
%\newtheorem{algorithm}{Algorithm}[section]
%\newtheorem{cor}{Corollary}
\newcommand{\BEQA}{\begin{eqnarray}}
\newcommand{\EEQA}{\end{eqnarray}}
\newcommand{\define}{\stackrel{\triangle}{=}}
\bibliographystyle{IEEEtran}
%\bibliographystyle{ieeetr}
\providecommand{\mbf}{\mathbf}
\providecommand{\pr}[1]{\ensuremath{\Pr\left(#1\right)}}
\providecommand{\qfunc}[1]{\ensuremath{Q\left(#1\right)}}
\providecommand{\sbrak}[1]{\ensuremath{{}\left[#1\right]}}
\providecommand{\lsbrak}[1]{\ensuremath{{}\left[#1\right.}}
\providecommand{\rsbrak}[1]{\ensuremath{{}\left.#1\right]}}
\providecommand{\brak}[1]{\ensuremath{\left(#1\right)}}
\providecommand{\lbrak}[1]{\ensuremath{\left(#1\right.}}
\providecommand{\rbrak}[1]{\ensuremath{\left.#1\right)}}
\providecommand{\cbrak}[1]{\ensuremath{\left\{#1\right\}}}
\providecommand{\lcbrak}[1]{\ensuremath{\left\{#1\right.}}
\providecommand{\rcbrak}[1]{\ensuremath{\left.#1\right\}}}
\makeatletter
\@addtoreset{figure}{section}
\makeatother
%\let\StandardTheFigure\thefigure
%\renewcommand{\thefigure}{\theproblem.\arabic{figure}}
\renewcommand{\thefigure}{\arabic{section}.\arabic{figure}}





\theoremstyle{remark}
\setcounter{figure}{0}
\newtheorem{rem}{Remark}
\newcommand{\sgn}{\mathop{\mathrm{sgn}}}
\providecommand{\abs}[1]{\left\vert#1\right\vert}
\providecommand{\res}[1]{\Res\displaylimits_{#1}} 
\providecommand{\norm}[1]{\left\lVert#1\right\rVert}
%\providecommand{\norm}[1]{\lVert#1\rVert}
\providecommand{\mtx}[1]{\mathbf{#1}}
\providecommand{\mean}[1]{E\left[ #1 \right]}
\providecommand{\fourier}{\overset{\mathcal{F}}{ \rightleftharpoons}}
%\providecommand{\hilbert}{\overset{\mathcal{H}}{ \rightleftharpoons}}
%\providecommand{\system}{\overset{\mathcal{H}}{ \longleftrightarrow}}
\providecommand{\system}[1]{\overset{\mathcal{#1}}{ \longleftrightarrow}}
\providecommand{\gauss}[2]{\mathcal{N}\ensuremath{\left(#1,#2\right)}}
%
	%\newcommand{\solution}[2]{\textbf{Solution:}{#1}}
\newcommand{\solution}{\noindent \textbf{Solution: }}
\newcommand{\cosec}{\,\text{cosec}\,}
\newcommand{\sinc}{\,\text{sinc}\,}
\newcommand{\rect}{\,\text{rect}\,}
\providecommand{\dec}[2]{\ensuremath{\overset{#1}{\underset{#2}{\gtrless}}}}
\newcommand{\myvec}[1]{\ensuremath{\begin{pmatrix}#1\end{pmatrix}}}
\newcommand{\mydet}[1]{\ensuremath{\begin{vmatrix}#1\end{vmatrix}}}
\newcommand*{\permcomb}[4][0mu]{{{}^{#3}\mkern#1#2_{#4}}}
\newcommand*{\perm}[1][-3mu]{\permcomb[#1]{P}}
\newcommand*{\comb}[1][-1mu]{\permcomb[#1]{C}}
%\numberwithin{equation}{section}
\numberwithin{equation}{section}
%\numberwithin{problem}{section}
%\numberwithin{definition}{section}
\makeatletter
\@addtoreset{figure}{problem}
\makeatother
\let\StandardTheFigure\thefigure
\let\vec\mathbf
\renewcommand{\thefigure}{\theproblem.\arabic{figure}}
\renewcommand{\thefigure}{\theproblem}
%\setlist[enumerate,1]{before=\renewcommand\theequation{\theenumi.\arabic{equation}}
%\counterwithin{equation}{enumi}
%\renewcommand{\theequation}{\arabic{subsection}.\arabic{equation}}
\vspace{3cm}
%\usepackage{babel}
\begin{document}
	
	
	
	\tikzstyle{block} = [rectangle, draw,
	text width=3em, text centered, minimum height=3em]
	\tikzstyle{sum} = [draw, circle, node distance=3cm]
	\tikzstyle{input} = [coordinate]
	\tikzstyle{output} = [coordinate]
	\tikzstyle{pinstyle} = [pin edge={to-,thin,black}]
	
	\theoremstyle{definition}

	%\newtheorem{algorithm}{Algorithm}[section]
	%\newtheorem{cor}{Corollary}

	\bibliographystyle{IEEEtran}
	%\bibliographystyle{ieeetr}
	\providecommand{\nCr}[2]{\,^{#1}C_{#2}} % nCr
	\providecommand{\nPr}[2]{\,^{#1}P_{#2}} % nPr
	\providecommand{\mbf}{\mathbf}
	\providecommand{\pr}[1]{\ensuremath{\Pr\left(#1\right)}}
	\providecommand{\qfunc}[1]{\ensuremath{Q\left(#1\right)}}
	\providecommand{\sbrak}[1]{\ensuremath{{}\left[#1\right]}}
	\providecommand{\lsbrak}[1]{\ensuremath{{}\left[#1\right.}}
	\providecommand{\rsbrak}[1]{\ensuremath{{}\left.#1\right]}}
	\providecommand{\brak}[1]{\ensuremath{\left(#1\right)}}
	\providecommand{\lbrak}[1]{\ensuremath{\left(#1\right.}}
	\providecommand{\rbrak}[1]{\ensuremath{\left.#1\right)}}
	\providecommand{\cbrak}[1]{\ensuremath{\left\{#1\right\}}}
	\providecommand{\lcbrak}[1]{\ensuremath{\left\{#1\right.}}
	\providecommand{\rcbrak}[1]{\ensuremath{\left.#1\right\}}}
	\theoremstyle{remark}
	
	\providecommand{\abs}[1]{\left\vert#1\right\vert}
	\providecommand{\res}[1]{\Res\displaylimits_{#1}} 
	\providecommand{\norm}[1]{\left\Vert#1\right\Vert}
	\providecommand{\mtx}[1]{\mathbf{#1}}
	\providecommand{\mean}[1]{E\left[ #1 \right]}
	\providecommand{\fourier}{\overset{\mathcal{F}}{ \rightleftharpoons}}
	%\providecommand{\hilbert}{\overset{\mathcal{H}}{ \rightleftharpoons}}
	\providecommand{\system}{\overset{\mathcal{H}}{ \longleftrightarrow}}
	%\newcommand{\solution}[2]{\textbf{Solution:}{#1}}
	
	\providecommand{\gauss}[2]{\mathcal{N}\ensuremath{\left(#1,#2\right)}}
	%
	\providecommand{\dec}[2]{\ensuremath{\overset{#1}{\underset{#2}{\gtrless}}}}
	\DeclarePairedDelimiter{\ceil}{\lceil}{\rceil}
	
	
	
	
\begin{tabular}{L{6cm} C{5cm} R{5cm} }
	\definecolor{circleorange}{rgb}{1,0.17,0.08}
\definecolor{darkorange}{rgb}{1,0.27,0.1}
\definecolor{orange2}{rgb}{1,0.5,0.15}
\definecolor{orange3}{rgb}{1,0.65,0.25}
\definecolor{yellow1}{rgb}{0.95,0.77,0.2}
%\begin{tikzpicture}[scale=0.2,every node/.style={transform shape}]
\begin{tikzpicture}[scale=0.1,every node/.style={transform shape}]
\draw [fill=circleorange,circleorange] (5,10) circle (1.15); 
\fill [darkorange] (5.06,8) -- (5.06,2) -- (7.3,1.2) -- (7.3,8.8) -- (5.06,8);
\fill [darkorange] (4.94,8) -- (4.94,2) -- (2.7,1.2) -- (2.7,8.8) -- (4.94,8);
\fill [orange2]    (7.4,8.4) -- (7.4,1.6) -- (8.2,1.2) -- (8.2,8.8) -- (7.4,8.4);
\fill [orange2]    (2.6,8.4) -- (2.6,1.6) -- (1.8,1.2) -- (1.8,8.8) -- (2.6,8.4);
\fill [orange3]    (8.3,8.4) -- (8.3,1.6) -- (9.0,1.2) -- (9.0,8.8) -- (8.3,8.4);
\fill [orange3]    (1.7,8.4) -- (1.7,1.6) -- (1.0,1.2) -- (1.0,8.8) -- (1.7,8.4);
\fill [yellow1]    (9.1,8.4) -- (9.1,1.6) -- (9.7,1.2) -- (9.7,8.8) -- (9.1,8.4);
\fill [yellow1]    (0.9,8.4) -- (0.9,1.6) -- (0.3,1.2) -- (0.3,8.8) -- (0.9,8.4);
\ifxetex
\node [scale=2.1] at (5,-0.1)  {   {\bf {\nakulafont  भारतीय प्रौद्योगिकी संस्थान हैदराबाद }} };
\node [scale=1.8] at (5,-1.2) {   {\bf { Indian Institute of Technology Hyderabad}} };
%\node [scale=1.8] at (5,-1.2) {   {\bf {\liberationsansfont Indian Institute of Technology Hyderabad}} };
\fi
\end{tikzpicture}
% \includegraphics[scale=0.05]{logo_iith} \newline
& Assignments 
	&
	AI1110
\end{tabular}
\vspace{-6mm}
\begin{center}
%\includegraphics[scale=0.95]{Yellow-Line}
\begin{tikzpicture}
\definecolor{yellow1}{rgb}{0.95,0.77,0.2}
\draw[line width=0.75mm, yellow1] (0,0) -- (\textwidth,0);
\end{tikzpicture}
\par\end{center}

\section{Uniform Random Numbers}
Let $U$ be a uniform random variable between 0 and 1.
\begin{enumerate}[label=\thesection.\arabic*
,ref=\thesection.\theenumi]
\item Generate $10^6$ samples of $U$ using a C program and save into a file called uni.dat .
\\
\solution Download the following files and execute the  C program.
\begin{lstlisting}
wget https://github.com/himanshukumargupta11012/Random-Numbers/blob/master/ques_1/1.1.c
wget https://github.com/himanshukumargupta11012/Random-Numbers/blob/master/coeffs.h
\end{lstlisting}
%
\item
Load the uni.dat file into python and plot the empirical CDF of $U$ using the samples in uni.dat. The CDF is defined as
\begin{align}
F_{U}(x) = \pr{U \le x}
\end{align}
\solution The CDF of U is plotted in Fig. \ref{fig:uni_cdf}
using the below python code
\begin{lstlisting}
wget https://github.com/himanshukumargupta11012/Random-Numbers/blob/master/ques_1/1.2.py
\end{lstlisting}
%\begin{figuree%%e%%\centeringe%\includegraphics[width=\umnwidth]{./ques_1/uni_cdf.png}
%\caption{The CDF of $U$}
%\label{fig:uni_cdf}
%e%%%%%
%
\item
Find a  theoretical expression for $F_{U}(x)$.
\\
\\
\\
\\
\\

\solution 
The CDF of U is given by
\begin{align}
	F_U(x)&=\pr{ U\leq x}\\
	&=\int_{-\infty}^x f_U(x)dx\\
\end{align}
Now, for $x<0$
\begin{align}
	f_U(x)=0\\
\therefore F_U(x)=0
\end{align}


,for $0\leq x \leq1$
\begin{align}
	f_U(x)&=1\\
\therefore F_U(x)&=\int_{-\infty}^x dx\\
	&=0+\int_{0}^x dx\\
	&=x
\end{align}
and for $x>1$
\begin{align}
	f_U(x)&=0\\
\therefore F_U(x)&=\int_{-\infty}^x dx\\
	&=0+\int_{0}^1 dx+0\\
	&=1
\end{align}

\[
So,\\
F_U(x)=
\begin{cases}
0&  x<0\\
 x  &  0 \leq x \leq 1\\
1&  x>1	
\end{cases}
\]

\begin{figure}
	\includegraphics[width=\columnwidth]{./ques_1/uni_cdf.png}
	\caption{The CDF of $U$}
	\label{fig:uni_cdf}
\end{figure}
\item
The mean of $U$ is defined as
%
\begin{equation}
E\sbrak{U} = \frac{1}{N}\sum_{i=1}^{N}U_i
\end{equation}
%
and its variance as
%
\begin{equation}
\text{var}\sbrak{U} = E\sbrak{U- E\sbrak{U}}^2 
\end{equation}
Write a C program to  find the mean and variance of $U$. 

\solution Download and run the following C program.
\begin{lstlisting}
wget https://github.com/himanshukumargupta11012/Random-Numbers/blob/master/ques_1/1.4.c
\end{lstlisting}


\item Verify your result theoretically given that
\begin{equation}
	E\sbrak{U^k} = \int_{-\infty}^{\infty}x^kdF_{U}(x)
\end{equation}
\solution The mean of U is given by
\begin{align}
	E[U]&=\int_{-\infty}^{\infty}xdF_U(x)\\
	&=0+\int_{0}^{1}xdx+0\\
	&=\bigg[\frac{x^2}{2}\bigg]_0^1\\
	&=1/2
\end{align}
and variance is given by,
\begin{align}
	E[U^2]-E[U]^2&=\int_{-\infty}^{\infty}x^2dF_U(x)-\brak{\frac{1}{2}}^2\\
	&=0+\int_{0}^{1}x^2dx+0-\brak{\frac{1}{2}}^2\\
	&=\Bigg[\frac{x^2}{3}\Bigg]_0^1-\brak{\frac{1}{2}}^2\\
	&=1/12
\end{align}
\end{enumerate}
%

\section{Central Limit Theorem}
%
\begin{enumerate}[label=\thesection.\arabic*
,ref=\thesection.\theenumi]
%
\item
Generate $10^6$ samples of the random variable
%
\begin{equation}
X = \sum_{i=1}^{12}U_i -6
\end{equation}
%
using a C program, where $U_i, i = 1,2,\dots, 12$ are  a set of independent uniform random variables between 0 and 1
and save in a file called gau.dat
%
\solution Download the following files and run the C program.
\begin{lstlisting}
wget https://github.com/himanshukumargupta11012/Random-Numbers/blob/master/ques_2/2.1.c
\end{lstlisting}
\item
Load gau.dat in python and plot the empirical CDF of $X$ using the samples in gau.dat. What properties does a CDF have?
\\
\solution The CDF of $X$ is plotted in Fig. \ref{fig:gauss_cdf}
using the below python code
\begin{lstlisting}
wget https://github.com/himanshukumargupta11012/Random-Numbers/blob/master/ques_2/2.2.py
\end{lstlisting}
Properties of CDF:
\begin{enumerate}
	\item It is always continuous.
	\item It is monotonic increasing.
	\item It is bounded between 0 to 1.
	\item At infinity, it tends to 1.
	\item At negative infinity, it tends to 0.
	\item Since it is bounded and have limit so it is convergent also.
	\item Q function is defined as:
	\begin{align}
		Q_X(x)=\pr{x<X}\\
	\therefore F_X\brak{x}=1-Q_X\brak{x}
	\end{align}
\end{enumerate}

\item
Load gau.dat in python and plot the empirical PDF of $X$ using the samples in gau.dat. The PDF of $X$ is defined as
\begin{align}
p_{X}(x) = \frac{d}{dx}F_{X}(x)
\end{align}
What properties does the PDF have?

\solution The PDF of $X$ is plotted in Fig. \ref{fig:gauss_pdf} using the below python code
\begin{lstlisting}
wget https://github.com/himanshukumargupta11012/Random-Numbers/blob/master/ques_2/2.2.py
\end{lstlisting}
Properties of PDF:
\begin{enumerate}
	\item It is always non-negative
	\item It is always bounded between 0 to 1
	\item Always symmetric about the mean.
	\item Maxima is observed at mean
	\item Area under the curve of PDF is always equal to 1
	\end{enumerate}
\begin{figure}
	\centering
	\includegraphics[width=\columnwidth]{./ques_2/gau_cdf.png}
	\caption{The CDF of $X$}
	\label{fig:gauss_cdf}
\end{figure}
\begin{figure}
\centering
\includegraphics[width=\columnwidth]{./ques_2/gau_pdf.png}
\caption{The PDF of $X$}
\label{fig:gauss_pdf}
\end{figure}
\item Find the mean and variance of $X$ by writing a C program.

\solution download the following files and run the C program.
\begin{lstlisting}
wget https://github.com/himanshukumargupta11012/Random-Numbers/blob/master/ques_2/2.4.c
wget https://github.com/himanshukumargupta11012/Random-Numbers/blob/master/coeffs.h
\end{lstlisting}
\item Given that 
\begin{align}
p_{X}(x) = \frac{1}{\sqrt{2\pi}}\exp\brak{-\frac{x^2}{2}}, -\infty < x < \infty,
\end{align}
repeat the above exercise theoretically.
%
\solution The mean of X is given by
\begin{align}
	E[X]&=\int_{-\infty}^\infty xp_X(x)dx\\
	&=\int_{-\infty}^\infty x\frac{1}{\sqrt{2\pi}}\exp\brak{-\frac{x^2}{2}}dx\\
	&=\Biggl[-\frac{1}{\sqrt{2\pi}}exp\brak{-\frac{x^2}{2}}\Biggl]_{-\infty}^{+\infty}
	&=0
\end{align}
and variance is given by
\begin{align}
	&E[X^2]-E[X]^2=\int_{-\infty}^\infty x^2f_X(x)dx\\
	&=\int_{-\infty}^\infty x^2\frac{1}{\sqrt{2\pi}}\exp\brak{-\frac{x^2}{2}}dx\\
	&=2\int_{0}^\infty x^2\frac{1}{\sqrt{2\pi}}\exp\brak{-\frac{x^2}{2}}dx\\
	&=\frac{2}{\sqrt{\pi}}\int_{0}^\infty t^{-1/2}\exp\brak{-t}dt\\
	&=\frac{2}{\sqrt{\pi}}\Gamma\brak{1+1/2}\\
	&=1
\end{align}

\end{enumerate}
\section{From Uniform to Other}
\begin{enumerate}[label=\thesection.\arabic*
,ref=\thesection.\theenumi]
%
\item
Generate samples of 
%
\begin{equation}
V = -2\ln\brak{1-U}
\end{equation}
%
and plot its CDF.  

\solution Download and run the python and C file where you get samples of V and fig. \ref{ray_cdf}
\begin{lstlisting}
wget https://github.com/himanshukumargupta11012/Random-Numbers/blob/master/ques_3/3.1.py
wget https://github.com/himanshukumargupta11012/Random-Numbers/blob/master/ques_3/3.1.c
\end{lstlisting}
\begin{figure}
	\centering
	\includegraphics[width=\columnwidth]{./ques_3/ray_cdf.png}
	\caption{The CDF of $V$}
	\label{fig:1}
\end{figure}
\item Find a theoretical expression for $F_V(x)$.

\solution The CDF of V is given by 
\begin{align}
	F_V(x)&=\pr{V\leq x}\\
	&=\pr{-2\ln\brak{1-U} \leq x}\\
	&=\pr{U\leq 1-\exp\brak{-\frac{x}{2}}}\\
	&=F_U\brak{1-\exp\brak{-\frac{x}{2}}}
\end{align}
let $a=\exp\brak{-x/2}$ .So,
\[
F_V(x)=
\begin{cases}
	0&  1-a<0\\
	1-a  &  0 \leq 1-a < 1\\
	1&  1-a\geq1	
\end{cases}
\]
on solving, we get
\[
F_V(x)=
\begin{cases}
	0&  x<0\\
	1-\exp\brak{-\frac{x}{2}}  &  x\geq0
\end{cases}
\]
%\item
%Generate the Rayleigh distribution from Uniform. Verify your result through graphical plots.
\end{enumerate}
\section{Triangular Distribution}
\begin{enumerate}[label=\thesection.\arabic*
,ref=\thesection.\theenumi]
%
\item Generate 
	\begin{align}
		T = U_1+U_2
	\end{align}

\solution Download the following code and run the C file
\begin{lstlisting}
wget https://github.com/himanshukumargupta11012/Random-Numbers/blob/master/ques_4/4.1.c

wget https://github.com/himanshukumargupta11012/Random-Numbers/blob/master/coeffs.h	
\end{lstlisting}
\item Find the CDF of $T$.

\solution
The CDF of T is given by
\begin{align}
	F_T\brak{t}=\pr{T\leq t}
\end{align}
Since T is sum of 2 uniform random variable between 0 and 1 
So,any T will be always less than 2 and greater than 0
So,
\begin{align}
	F_T\brak{t}=0 \hspace{10mm} \text{for $X<0$}\\
	F_T\brak{t}=1 \hspace{10mm} \text{for $X\geq0$}
\end{align}
Now,for $0\leq t<2$,

let $U_1$ as a constant x and later we integrate it for all possible values of x
So,
\begin{align}
	F_T\brak{t}&=\pr{T\leq t}\\
	&=\pr{U_1+U_2\leq t}\\
&=\pr{U_2\leq t-U_1}\\
&=\int_0^t F_{U_2}\brak{t-x} p_{U_1}\brak{x}dx
\end{align}
Now, we make 2 cases


\textbf{case-1:} $0\leq t<1$
\begin{align}
	F_T\brak{t}&=\int_0^t F_{U_2}\brak{t-x}dx\\
	&=\int_0^t F_{U_2}\brak{m}dm\\
	&=\int_0^t mdm\\
	&=\bigg[\frac{m^2}{2}\bigg]_0^t\\
	&=\frac{t^2}{2}
\end{align}
\textbf{case-2:} $1\leq t<2$
\begin{align}
	F_T\brak{t}&=\int_0^1 F_{U_2}\brak{t-x}dx\\
	&=\int_0^1 F_{U_2}\brak{t-x}dx\\
	&=\int_{t-1}^t F_{U_2}\brak{m}dm\\
	&=\int_{t-1}^1 mdm+\int_1^t dm\\
	&=\bigg[\frac{m^2}{2}\bigg]_{t-1}^1+[1]_1^t\\
	&=\frac{-t^2}{2}+2t-1	
\end{align}

\[
\therefore
F_T(t)=
\begin{cases}
	0&  t<0\\
	\frac{t_2}{2} &  0\leq t<1\\
	\frac{-t^2}{2}+2t-1  &  1\leq t<2\\
	1  & 2\leq t
\end{cases}
\]
\item Find the PDF of $T$.

\solution The PDF of T is given by
\begin{align}
	p_{T}(t) = \frac{d}{dt}F_{T}(t)
\end{align}
\[
\therefore
F_T(t)=
\begin{cases}
	0&  t<0\\
	t &  0\leq t<1\\
	2-t  &  1\leq t<2\\
	0  & 2\leq t
\end{cases}
\]
\item Find the theoretical expressions for the PDF and CDF of $T$.

\solution Already done in 4.2 and 4.3.
\item Verify your results through a plot. 

\solution Plot of CDF and PDF is plotted in fig. 

\begin{figure}
	\centering
		\includegraphics[width=\columnwidth]{./ques_4/tri_cdf.png}
	\caption{The CDF of $T$}
	\label{tri_cdf}
\end{figure}

\begin{figure}
	\centering
		\includegraphics[width=\columnwidth]{./ques_4/tri_pdf.png}
	\caption{The PDF of $T$}
	\label{tri_pdf}
\end{figure}
\eqref{tri_cdf} and \eqref{tri_pdf} respectively using following python code.
\end{enumerate}


\begin{comment}
\item Generate 
\begin{equation}
Y = AX+N,
\end{equation}
		where $A = 5 \text{ dB}, X \i \cbrak{1,-1}$,  is Bernoulli and $N \sim \gauss{0}{1}$.
		
	\solution Download the following files and execute the C code
	\begin{lstlisting}
wget https://github.com/himanshukumargupta11012/Random-Numbers/blob/master/ques_5/5.1.c
wget https://github.com/himanshukumargupta11012/Random-Numbers/blob/master/coeffs.h
	\end{lstlisting}
	\item Plot $Y$.
	
	\solution
	The plot of Y is plotted in fig. using below python code
	\begin{lstlisting}
wget https://github.com/himanshukumargupta11012/Random-Numbers/blob/master/ques_5/5.2.py
	\end{lstlisting} 
	\item Guess how to estimate $X$ from $Y$.
\item
\label{ml-ch4_sim}
Find 
\begin{equation}
	P_{e|0} = \pr{\hat{X} = -1|X=1}
\end{equation}
and 
\begin{equation}
	P_{e|1} = \pr{\hat{X} = 1|X=-1}
\end{equation}
%
\item Find $P_e$.
%
\item
Verify by plotting  the theoretical $P_e$.  
		\end{enumerate}
	
\section{Gaussian to Other}

	

\begin{enumerate}[label=\thesection.\arabic*
,ref=\thesection.\theenumi]
\item
Let $X_1 \sim  \gauss{0}{1}$ and $X_2 \sim  \gauss{0}{1}$. Plot the CDF and PDF of
%
\begin{equation}
V = X_1^2 + X_2^2
\end{equation}

\solution 
Download the following files and run C and python program.
CDF and PDF of V is plotted in fig. \ref{} and \ref{} respectively using below python code.
\begin{lstlisting}
wget https://github.com/himanshukumargupta11012/Random-Numbers/blob/master/coeffs.h
wget https://github.com/himanshukumargupta11012/Random-Numbers/blob/master/ques_6/6.1.c
wget https://github.com/himanshukumargupta11012/Random-Numbers/blob/master/ques_6/6.1.py
\end{lstlisting}

\begin{figure}
	\centering
	\includegraphics[width=\columnwidth]{./ques_6/gau_gau_cdf.png}
	\caption{The DF of $V$}
	\label{gau_gau_cdf}
\end{figure}
\begin{figure}
	\centering
	\includegraphics[width=\columnwidth]{./ques_6/gau_gau_pdf.png}
	\caption{The PDF of $V$}
	\label{gau_gau_pdf}
\end{figure}
\item
If
%
\begin{equation}
F_{V}(x) = 
\begin{cases}
1 - e^{-\alpha x} & x \geq 0 \\
0 & x < 0,
\end{cases}
\end{equation}
%
find $\alpha$.

\solution Let 

$X_1=R\cos\theta$ ,$X_2=R\sin\theta$
So,

\item
\label{ch3_raleigh_sim}
Plot the CDF and PDf of
%
\begin{equation}
A = \sqrt{V}
\end{equation}
\solution
CDF and PDF of A is plotted in fig. \ref{} and \ref{} respectively using below python code.
\begin{lstlisting}
wget https://github.com/himanshukumargupta11012/Random-Numbers/blob/master/ques_6/6.3.py
	\end{lstlisting}
\end{enumerate}
\section{Conditional Probability}
\begin{enumerate}[label=\thesection.\arabic*
,ref=\thesection.\theenumi]
\item
\item
\label{ch4_sim}
Plot 
\begin{equation}
P_e = \pr{\hat{X} = -1|X=1}
\end{equation}
%
for 
\begin{equation}
Y = AX+N,
\end{equation}
where $A$ is Raleigh with $E\sbrak{A^2} = \gamma, N \sim \gauss{0}{1}, X \in \brak{-1,1}$ for $0 \le \gamma \le 10$ dB.
%
\item
Assuming that $N$ is a constant, find an expression for $P_e$.  Call this $P_e(N)$
%
\item
%
\label{ch4_anal}
For a function $g$,
\begin{equation}
E\sbrak{g(X)} = \int_{-\infty}^{\infty}g(x)p_{X}(x)\, dx
\end{equation}
%
Find $P_e = E\sbrak{P_e(N)}$.
%
\item
Plot $P_e$ in problems \ref{ch4_sim} and \ref{ch4_anal} on the same graph w.r.t $\gamma$.  Comment.
		\end{enumerate}
\section{Two Dimensions}
Let 
\begin{equation}
\mbf{y} = A\mbf{x} + \mbf{n},
\end{equation}
where 
\begin{align}
x &\in \brak{\mbf{s}_0,\mbf{s}_1}, 
\mbf{s}_0 = 
\begin{pmatrix}
1 
\\
0
\end{pmatrix},
\mbf{s}_1 = 
\begin{pmatrix}
0 
\\
1
\end{pmatrix}
\\
\mbf{n} &= 
\begin{pmatrix}
n_1
\\
n_2
\end{pmatrix},
n_1,n_2 \sim \gauss{0}{1}.
\end{align}
%
\begin{enumerate}[label=\thesection.\arabic*
,ref=\thesection.\theenumi]
%%
\item
\label{ch5_fsk}
Plot 
%
\begin{equation}
\mbf{y}|\mbf{s}_0 \text{ and } \mbf{y}|\mbf{s}_1
\end{equation}
%
on the same graph using a scatter plot.
%
\item
For the above problem, find a decision rule for detecting the symbols $\mbf{s}_0 $ and $\mbf{s}_1$.
%
\item
Plot 
\begin{equation} 
P_e = \pr{\hat{\mbf{x}} = \mbf{s}_1|\mbf{x} = \mbf{s}_0}
\end{equation}
with respect to the SNR from 0 to 10 dB.
%
\item
Obtain an expression for $P_e$. Verify this by comparing the theory and simulation plots on the same graph.
%
		\end{enumerate}
\end{document}
\end{comment}
\section{Maximul Likelihood}

\begin{enumerate}[label=\thesection.\arabic*
	,ref=\thesection.\theenumi]

\item Generate equiprobable $X \in \cbrak{1,-1}$.

\solution Download the following files and execute the C file
\begin{lstlisting}
wget https://github.com/himanshukumargupta11012/Random-Numbers/blob/master/ques_5/5.1.c
wget https://github.com/himanshukumargupta11012/Random-Numbers/blob/master/coeffs.h
\end{lstlisting}
\item Generate 
\begin{equation}
	Y = AX+N,
\end{equation}
where $A = 5$ dB,  and $N \sim \gauss{0}{1}$.
\solution Download the following files and run the C code
\begin{lstlisting}
wget https://github.com/himanshukumargupta11012/Random-Numbers/blob/master/ques_5/5.2.c
wget https://github.com/himanshukumargupta11012/Random-Numbers/blob/master/coeffs.h
\end{lstlisting}
\item Plot $Y$ using a scatter plot.

\solution The scatter plot of Y is plotted in fig. \ref{scatter_plot} using below python code.
\begin{lstlisting}
wget https://github.com/himanshukumargupta11012/Random-Numbers/blob/master/ques_5/5.3.py
\end{lstlisting}
\begin{figure}
	\centering
	\includegraphics[width=\columnwidth]{./ques_5/scatter_plot.png}
	\caption{The scatter plot of $Y$}
	\label{scatter_plot}
\end{figure}
\item Guess how to estimate $X$ from $Y$.

\solution After seeing the scatter plot we can conclude that  if Y is positive then X will be 1 and if Y is negative then X will be -1.This conclusion also depends on value of A because if we decrease A the two parts of graph intermix and then we can't say anything.If A is approximately greater than 3 then only we can conclude this.
\item
\label{ml-ch4_sim}
Find 
\begin{equation}
	P_{e|0} = \pr{\hat{X} = -1|X=1}
\end{equation}
and 
\begin{equation}
	P_{e|1} = \pr{\hat{X} = 1|X=-1}
\end{equation}
%
\solution $\hat{X}$ is defined as


\begin{align}
	\hat{X}=
\begin{cases}
	1  &  Y\geq 0\\
	-1 &  Y<0
\end{cases}
\end{align}
So,
\begin{align}
	P_{e|0}&=\pr{\hat{X}=-1|X=1}\\
	&=\pr{Y<0|X=1}\\
	&=\pr{A\times 1+N<0}\\
	&=\pr{N<-A}\\
	&=F_N\brak{-A}\\
	&=0
\end{align}
%	&=\int_{-\infty}^{-A}p_N\brak{x}dx\\
%	&=\int_{-\infty}^{-A}\frac{1}{\sqrt{2\pi}}\exp\brak{%-\frac{x^2}{2}}dx
and 
\begin{align}
	P_{e|1}&=\pr{\hat{X}=1|X=-1}\\
	&=\pr{Y\geq0|X=-1}\\
	&=\pr{A\times -1+N\geq0}\\
	&=\pr{N\geq A}\\
	&=1-\pr{N<A}\\
	&=1-F_N\brak{A}\\
	&=0
\end{align}
Download and run the following file to get coding proof
\begin{lstlisting}
wget https://github.com/himanshukumargupta11012/Random-Numbers/blob/master/ques_5/5.5.py
\end{lstlisting}
\item Find $P_e$ assuming that $X$ has equiprobable symbols.

\solution  $P_e$ can be defined as 
\begin{align}
	P_e&=\pr{X=-1}P_{e|1}+\pr{X=1}P_{e|0}\\
	&\because P_{e|1}=0,P_{e|0}=0\\
	&\therefore P_e=0
\end{align}

\item
Verify by plotting  the theoretical $P_e$ with respect to $A$ from 0 to 10 dB.  

\solution X is equiprobable.So
\begin{align}
	\because& \pr{X=1}=\pr{X=-1},\\
	&\pr{X=1}+\pr{X=-1}=1\\
\therefore&	\pr{X=1}=\pr{X=-1}=1/2
\end{align}
Now,$P_e$ is given by
\begin{align}
	P_e&=\pr{X=-1}P_{e|1}+\pr{X=1}P_{e|0}\\
	P_e&=\frac{1}{2}\brak{1-F_N\brak{A}}+\frac{1}{2}\brak{F_N\brak{-A}}
\end{align}
Now, the plot of $P_e$ w.r.t. A is in fig. \ref{p_e_the} by plotting the above function using following code
\begin{lstlisting}
wget https://github.com/himanshukumargupta11012/Random-Numbers/blob/master/ques_5/5.7.py
\end{lstlisting} 

\begin{figure}
\centering
\includegraphics[width=\columnwidth]{./ques_5/p_e_the.png}
\caption{$P_e$ graph wrt A}
\label{p_e_the}
\end{figure}
\item Now, consider a threshold $\delta$  while estimating $X$ from $Y$. Find the value of $\delta$ that maximizes the theoretical $P_e$.

\solution Taking threshold as $\delta$ we get
\begin{align}
	P_{e|0}&=\pr{Y<\delta| X=1}\\
	&=\pr{A\times 1+N<\delta}\\
	&=\pr{N<-A+\delta}\\
	&=F_N\brak{-A+\delta}
\end{align}
and \begin{align}
	P_{e|1}&=\pr{Y\geq \delta|X=-1}\\
	&=\pr{A\times -1+N\geq \delta}\\
	&=\pr{N\geq A+\delta}\\
	&=1-\pr{N<A+\delta}\\
	&=1-F_N\brak{A+\delta}
\end{align}
and $P_e$ is defined as 
\begin{align}
	P_e &=\pr{X=-1}P_{e|1}+\pr{X=1}P_{e|0}\\
	&=\frac{1}{2}\brak{1-F_N\brak{A+\delta}+F_N\brak{-A+\delta}}
\end{align}
Now we differentiate $P_e$ and equate it to 0 to get $\delta$  that minimizes the $P_e$.So,
\begin{align}
&\frac{d}{d\delta}P_e=\frac{1}{2}\brak{\frac{d}{d\delta} F_N\brak{-A+\delta}-\frac{d}{d\delta} F_N\brak{A+\delta}}=0\\
&\implies p_N\brak{-A+\delta}-p_N\brak{A+\delta}=0\\
&\implies \exp\brak{-\frac{\brak{-A+\delta}^2}{2}}-\exp\brak{-\frac{\brak{A+\delta}^2}{2}}=0\\
&\implies \delta=0
\end{align}

Now we have to check that $\delta$=0 is for maxima or minima by double differentiating it.
So,
\begin{align}
	\frac{d^2}{d\delta^2}|_{\delta=0}&=2A\exp\brak{\frac{-A^2}{2}}\\
	&=10\exp\brak{\frac{-A^2}{2}}
	&>0
\end{align}
Double derivative is positive.
So,$\delta$=0 will give minima of $P_e$

Plot of $P_e$ w.r.t. $\delta$ is plotted in fig. \ref{p_e_del} using following code
\begin{lstlisting}
wget https://github.com/himanshukumargupta11012/Random-Numbers/blob/master/ques_5/5.8.py
\end{lstlisting}
\begin{figure}
	\centering
	\includegraphics[width=\columnwidth]{./ques_5/p_e_del.png}
	\caption{$P_e$ graph wrt $\delta$}
	\label{p_e_del}
\end{figure}
%&=\frac{1}{2}\brak{\frac{erfc(\brak{A+\delta}/\sqrt{2})}{2}}\\
%&+\frac{1}{2}\brak{-\frac{erfc\brak{\brak{-A+\delta}/\sqrt{2}}}{2}+1}
\item Repeat the above exercise when 
\begin{align}
	p_{X}(0) = p
\end{align}
\solution 
\begin{align}
\because	\pr{X=1}&=p_X\brak{0}=p\\
\therefore\pr{X=-1}&=1-\pr{X=1}\\
&=1-p
\end{align}
and $P_e$ is given by 
\begin{align}
	P_e &=\pr{X=-1}P_{e|1}+\pr{X=1}P_{e|0}\\
	&=\brak{1-p}\brak{1-F_N\brak{A+\delta}}\\
	&+pF_N\brak{-A+\delta}
\end{align}
Then,
\begin{align}
	&\frac{d}{d\delta}P_e=p\frac{d}{d\delta} F_N\brak{-A+\delta}\\
	&-\brak{1-p}\frac{d}{d\delta} F_N\brak{A+\delta}=0\\
	&\implies p\times p_N\brak{-A+\delta}\\
	&-\brak{1-p}p_N\brak{A+\delta}=0\\
	&\implies p\times \exp\brak{-\frac{\brak{-A+\delta}^2}{2}}
	\\&-\brak{1-p}\exp\brak{-\frac{\brak{A+\delta}^2}{2}}=0\\
	&\implies \delta=\frac{\log\brak{\frac{1}{p}-1}}{2A}
\end{align}
\item Repeat the above exercise using the MAP criterion.

\solution From bayes theorem we can write,
\begin{align}
&\pr{X=1|Y=y}\\
&=\frac{\pr{Y=y|X=1}\pr{X=1}}{\pr{Y=y}}\\
&=\frac{\pr{N=y-A}p}{\pr{N=y-A}p+\pr{N=y+A}\brak{1-p}}\\
&=\frac{p_N\brak{y-A}p}{p_N\brak{y-A}p+p_N\brak{y+A}\brak{1-p}}\\
&=\frac{p}{p+\brak{1-p}\exp\brak{-2Ay}}
\end{align}
%\frac{\pr{A\times1+N=0}p}

	\begin{align}
&\pr{X=-1|Y=y}\\
&=1-\pr{X=1|Y=y}\\
&=1-\frac{p}{p+\brak{1-p}\exp\brak{-2Ay}}\\
&=\frac{\brak{1-p}\exp\brak{-2Ay}}{p+\brak{1-p}\exp\brak{-2Ay}}\\
&=\frac{1-p}{p\exp\brak{2Ay}+1-p}
	\end{align}
Now, if $X=1$ is more likely than $X=-1$ for a y. Then,
\begin{align}
	\pr{X=1|Y=y}&>\pr{X=-1|Y=y}\\
	\frac{p}{p+\brak{1-p}\exp\brak{-2Ay}}&>\frac{1-p}{p\exp\brak{2Ay}+1-p}\\
	y>\frac{\log{\brak{\frac{1-p}{p}}}}{2A}
\end{align} 
and if $X=-1$ is more lileky for a y. Then,
\begin{align}
	&\pr{X=1|Y=y}<\pr{X=-1|Y=y}\\
	&\frac{p}{p+\brak{1-p}\exp\brak{-2Ay}}\\
	&<\frac{1-p}{p\exp\brak{2Ay}+1-p}\\
	&y<\frac{\log{\brak{\frac{1-p}{p}}}}{2A}
\end{align}
So,we can see that optimal value of Y for $P_e$ to be minimum is 
\begin{align}
	\frac{\log{\brak{\frac{1-p}{p}}}}{2A}
\end{align}
So,
\begin{align}
	\delta=\frac{\log{\brak{\frac{1-p}{p}}}}{2A}
\end{align}
\end{enumerate}
\section{Gaussian to Other}
\begin{enumerate}[label=\thesection.\arabic*
,ref=\thesection.\theenumi]
\item
Let $X_1 \sim  \gauss{0}{1}$ and $X_2 \sim  \gauss{0}{1}$. Plot the CDF and PDF of
%
\begin{equation}
	V = X_1^2 + X_2^2
\end{equation}
\solution 
Download the following files and run C and python program.
CDF and PDF of V is plotted in fig. \ref{gau_gau_cdf} and \ref{gau_gau_pdf} respectively using below python code.
\begin{lstlisting}
wget https://github.com/himanshukumargupta11012/Random-Numbers/blob/master/coeffs.h
wget https://github.com/himanshukumargupta11012/Random-Numbers/blob/master/ques_6/6.1.c
wget https://github.com/himanshukumargupta11012/Random-Numbers/blob/master/ques_6/6.1.py
\end{lstlisting}

\begin{figure}
	\centering
	\includegraphics[width=\columnwidth]{./ques_6/gau_gau_cdf.png}
	\caption{The CDF of $V$}
	\label{gau_gau_cdf}
\end{figure}
\begin{figure}
	\centering
	\includegraphics[width=\columnwidth]{./ques_6/gau_gau_pdf.png}
	\caption{The PDF of $V$}
	\label{gau_gau_pdf}
\end{figure}
\item
If
%
\begin{equation}
	F_{V}(x) = 
	\begin{cases}
		1 - e^{-\alpha x} & x \geq 0 \\
		0 & x < 0,
	\end{cases}
\end{equation}
%
find $\alpha$.

\solution Let 

$X_1=R\cos\theta$ ,$X_2=R\sin\theta$

So,
PDF of R and $\theta$ is given by 
\begin{align}
	p_{R,\theta}\brak{r,\phi}=p_{X_1,X_2}\brak{x_1,x_2}|J|
\end{align} 
where J is jacobian matrix transforming R,$\theta$ to $X_1,X_2$

On soving we get jacobian$,J=R.$
So,
\begin{align}
	p_{R,\theta}\brak{r,\phi}=p_{X_1,X_2}\brak{x_1,x_2}r
\end{align} 
since $X_1$ and $X_2$ are i.i.d. and gaussian random variable,$\gauss{0}{1}$.So,
\begin{align}
	p_{R,\theta}\brak{r,\phi}&=p_{X_1}\brak{x_1}p_{X_2}\brak{x_2}r\\
	&=\frac{1}{\sqrt{2\pi}}\exp\brak{-\frac{x_1^2}{2}}\frac{1}{\sqrt{2\pi}}\exp\brak{-\frac{x_2^2}{2}}r\\
	&=\frac{1}{2\pi}\exp\brak{-\frac{r^2}{2}}r
\end{align}
Now PDF of R,
\begin{align}
	p_R\brak{r}&=\int_0^{2\pi}p_{R,\theta}\brak{r,\phi}d\phi\\
	&=\int_0^{2\pi} \frac{1}{2\pi}\exp\brak{-\frac{r^2}{2}}rd\phi\\
	&=\exp\brak{-\frac{r^2}{2}}r
\end{align}
and
\begin{align}
	V=X_1^2 + X_2^2=R^2
\end{align}
Since $V=R^2>0$

So,CDF and PDF would be zero for v$<$0


%So,PDF of V,
%\begin{align}
%p_V\brak{v}&=\pr{V=v}\\
%&=\pr{R^2=v}\\
%&=\pr{R=\sqrt v}\\
%&=p_R\brak{\sqrt{v}}\\
%&=\exp\brak{-\frac{v}{2}}\sqrt v
%\end{align}
%Then CDF of V,
%\begin{align}
%F_V\brak{v}&=\int_{-\infty}^v p_V\brak{v}dv\\
%&=\int_{-\infty}^v \exp\brak{-\frac{v}{2}}\sqrt v dv
%\end{align}
Then CDF of V for v$\geq$0, 
\begin{align}
F_V\brak{v}&=\pr{V\leq v}\\
&=\pr{R^2\leq v}\\
&=\pr{R\leq \sqrt v}\\
&=F_R\brak{\sqrt v}\\
&=\int_{-\infty}^{\sqrt v} p_R\brak{r}dr\\
&=0+\int_{0}^{\sqrt v}\exp\brak{-\frac{r^2}{2}}r dr\\
&=1-\exp\brak{-\frac{v}{2}}
\end{align}
Comparing the 2 equations we get 
\begin{align}
	\alpha=\frac{1}{2}
\end{align}
and PDF is given by
\begin{align}
	p_V\brak{v}&=\frac{d}{dx}F_V\brak{v}\\
	\therefore p_V\brak{v}&=\begin{cases}
		\frac{\exp\brak{-\frac{v}{2}}}{2}& v\geq 0\\
		0  & v<0
	\end{cases}
\end{align}

\item
\label{ch3_raleigh_sim}
Plot the CDF and PDf of
%
\begin{equation}
	A = \sqrt{V}
\end{equation}
\solution
CDF and PDF of A is plotted in fig. \ref{gau_gau_rt_cdf} and \ref{gau_gau_rt_pdf} respectively using below python code.
\begin{lstlisting}
wget https://github.com/himanshukumargupta11012/Random-Numbers/blob/master/ques_6/6.3.py
\end{lstlisting}
We know that 
\begin{align}
	V&=R^2\\
\therefore A&= \sqrt V=R
\end{align}
Since $A=R>0$. So PDF and CDF will be 0 for $a<0$

Then,PDF of A for $a\geq 0$,
\begin{align}
	F_A\brak{a}&=\pr{A\leq a}\\
	&=\pr {R\leq a}\\
	&=F_R\brak{a}\\
&=\int_{-\infty}^{a} p_R\brak{r}dr\\
&=0+\int_{0}^{a}\exp\brak{-\frac{r^2}{2}}r dr\\
&=1-\exp\brak{-\frac{a^2}{2}}
\end{align}
So,CDF of A,
\begin{align}
	F_A\brak{a}=
\begin{cases}
	1-\exp\brak{-\frac{a^2}{2}} & a\geq 0\\
	0 &a<0
\end{cases}
\end{align}

PDF of A is given by
\begin{align}
p_A\brak{a}&=\frac{d}{dx}F_A\brak{a}\\
\therefore p_A\brak{a}&= \begin{cases}
			\exp\brak{-\frac{a^2}{2}}a & a\geq 0\\
			0 &a<0
		\end{cases}
\end{align}

\begin{figure}
	\centering
	\includegraphics[width=\columnwidth]{./ques_6/gau_gau_rt_cdf.png}
	\caption{The CDF of $A$}
	\label{gau_gau_rt_cdf}
\end{figure}
\begin{figure}
	\centering
	\includegraphics[width=\columnwidth]{./ques_6/gau_gau_rt_pdf.png}
	\caption{The PDF of $A$}
	\label{gau_gau_rt_pdf}
\end{figure}
%Then,PDF of A for $a\geq 0$,
%\begin{align}
%	p_A\brak{a}&=\pr{A=a}\\
%	&=\pr {R=a}\\
%	&=p_R\brak{a}\\
%	&=\exp\brak{-\frac{a^2}{2}}a
%\end{align}
%So,PDF of A,
%\begin{align}
%	p_A\brak{a}=
%	\begin{cases}
%		\exp\brak{-\frac{a^2}{2}}a & a\geq 0\\
%		0 &a<0
%	\end{cases}
%\end{align}
%We know that 
%\begin{align}
%	F_A\brak{a}=\int_{-\infty}^a p_A\brak{x} dx
%\end{align}
%So,CDF of A,
%\begin{align}
%	F_A\brak{a}=
%\begin{cases}
%	ddgdg & a\geq 0\\
%	0 &a<0
%\end{cases}
%\end{align}

\end{enumerate}
\section{Conditional Probability}
\begin{enumerate}[label=\thesection.\arabic*
,ref=\thesection.\theenumi]
\item
\label{ch4_sim}
Plot 
\begin{equation}
	P_e = \pr{\hat{X} = -1|X=1}
\end{equation}
%
for 
\begin{equation}
	Y = AX+N,
\end{equation}
where $A$ is Raleigh with $E\sbrak{A^2} = \gamma, N \sim \gauss{0}{1}, X \in \brak{-1,1}$ for $0 \le \gamma \le 10$ dB.
%
\item
Assuming that $N$ is a constant, find an expression for $P_e$.  Call this $P_e(N)$
%
\item
%
\label{ch4_anal}
For a function $g$,
\begin{equation}
	E\sbrak{g(X)} = \int_{-\infty}^{\infty}g(x)p_{X}(x)\, dx
\end{equation}
%
Find $P_e = E\sbrak{P_e(N)}$.
%
\item
Plot $P_e$ in problems \ref{ch4_sim} and \ref{ch4_anal} on the same graph w.r.t $\gamma$.  Comment.
\end{enumerate}
\section{Two Dimensions}
Let 
\begin{equation}
\mbf{y} = A\mbf{x} + \mbf{n},
\end{equation}
where 
\begin{align}
x &\in \brak{\mbf{s}_0,\mbf{s}_1}, 
\mbf{s}_0 = 
\begin{pmatrix}
	1 
	\\
	0
\end{pmatrix},
\mbf{s}_1 = 
\begin{pmatrix}
	0 
	\\
	1
\end{pmatrix}
\\
\mbf{n} &= 
\begin{pmatrix}
	n_1
	\\
	n_2
\end{pmatrix},
n_1,n_2 \sim \gauss{0}{1}.
\end{align}
%
\begin{enumerate}[label=\thesection.\arabic*
,ref=\thesection.\theenumi]
%%
\item
\label{ch5_fsk}
Plot 
%
\begin{equation}
	\mbf{y}|\mbf{s}_0 \text{ and } \mbf{y}|\mbf{s}_1
\end{equation}
%
on the same graph using a scatter plot.
%
\item
For the above problem, find a decision rule for detecting the symbols $\mbf{s}_0 $ and $\mbf{s}_1$.
%
\item
Plot 
\begin{equation} 
	P_e = \pr{\hat{\mbf{x}} = \mbf{s}_1|\mbf{x} = \mbf{s}_0}
\end{equation}
with respect to the SNR from 0 to 10 dB.
%
\item
Obtain an expression for $P_e$. Verify this by comparing the theory and simulation plots on the same graph.
%
\end{enumerate}
\end{document}